% =================================================================
\chapter{Logstruktur und Strategie}
% =================================================================

Je nach Kundenanforderungen hinsichtlich gesetzlicher Vorgaben, Aufbewahrungszeiträumen, Infrastruktur und verfügbaren Kapazitäten kann die Struktur, in der Logdateien abgelegt werden, variieren.

In der von LogObject empfohlenen Struktur werden alle Logdateien in einen zentralen Logordner geschrieben, sodass sämtliche Logs systemweit an einer Stelle einsehbar sind.

\section{Logstruktur und Rollover}

Für Überwachung und Analyse ist es in den meisten Umgebungen am einfachsten, wenn alle Logdateien in einem zentralen Logverzeichnis abgelegt werden.
Bei SNDE wird hierfür der Pfad \texttt{/local/log} verwendet.

Dort werden die Logs servicebezogen (z.\,B. \textit{mLogistics}, \textit{LoMFJ}, \textit{LoRouter} usw.) in separaten Unterverzeichnissen strukturiert.
Innerhalb jedes Service-Verzeichnisses wird für jeden Logstream ein eigener Unterordner angelegt, beispielsweise:

\begin{itemize}
	\item \texttt{log/lomfj/audit}
	\item \texttt{log/lomfj/mfjserver}
	\item \texttt{log/lorouter/server}
	\item \texttt{log/lorouter/error}
\end{itemize}

In diesen Unterordnern befindet sich jeweils die aktuell laufende Logdatei, in die fortlaufend Logeinträge geschrieben werden.

\subsection{Rollover}

Beim Log-Rollover wird die aktuelle Logdatei automatisch komprimiert (\texttt{.gz} oder \texttt{.zip}) und kann anschließend in ein beliebiges Zielverzeichnis verschoben werden.

Üblich ist entweder eine rein lokale Ablage oder eine strukturierte Ablage in tagesbezogenen Unterordnern mit Datumsangabe, beispielsweise:

\begin{itemize}
	\item \texttt{/log/lorouter/error/1970-01-01/log.1.gz}
	\item \texttt{/log/lorouter/error/1970-01-01/log.2.gz}
\end{itemize}

Der Zielpfad für diese Rollover-Archive ist frei konfigurierbar und kann beispielsweise auch auf einen Fileserver verweisen.

Vom direkten Schreiben der laufenden Logdateien auf einen Fileserver wird ausdrücklich abgeraten, da dies zu Performance- und Stabilitätsproblemen führen kann.
Für die Verwaltung und den Betrieb der auf dem Fileserver abgelegten Logarchive ist der jeweilige Fileserver-Betreiber verantwortlich.

\subsubsection{Kompression}

Logdateien lassen sich aufgrund ihres stark repetitiven, textbasierten Inhalts sehr effizient komprimieren. In der Praxis sind Platzersparnisse von 85--95\,\% gegenüber der unkomprimierten Datei üblich.

Für die Kompression beim Rollover ist das \texttt{.gz}-Format dem \texttt{.zip}-Format klar überlegen. \texttt{.gz} arbeitet als reiner Datenstrom und ist dadurch schneller, benötigt weniger CPU- und I/O-Ressourcen und erzielt in der Regel eine bessere Kompressionsrate. \texttt{.zip} hingegen ist ein Archivformat mit zusätzlichem Overhead, das vor allem für mehrere Dateien gedacht ist und bei einzelnen Logdateien unnötige Ressourcen verbraucht.

Daher gilt für Logdateien als Best Practice die Verwendung von \texttt{.gz} statt \texttt{.zip}.

\subsubsection{Löschen}

Unabhängig davon, wo die Rollover-Dateien gespeichert werden und ob diese in Datumsordnern gruppiert sind, kann eine zeitbasierte Löschstrategie konfiguriert werden.

\section{Lokale Konfiguration}

Die Logs werden nach folgenden Regeln geschrieben:

\begin{itemize}
	\item Laufende Logs
	      \begin{itemize}
		      \item werden lokal auf den Maschinen geführt,
		      \item werden im Pfad \texttt{/local/log/applicationName} gespeichert,
		      \item folgen dem Namensschema \texttt{\textit{applicationName}\_\textit{logName}}.
	      \end{itemize}
	\item Rollover
	      \begin{itemize}
		      \item \item erfolgen auf den Fileserver unter \path{/els/shared/{hostname}/applicationName/logName/dailyRollover}.
		      \item werden als \texttt{.gz} komprimiert.
		      \item Ein Rollover wird ausgelöst durch:
		            \begin{itemize}
			            \item Erreichen einer Größe von 20\,MB,
			            \item Neustart der Applikation,
			            \item täglich um Mitternacht.
		            \end{itemize}
	      \end{itemize}
	\item Logdateien werden nach einer konfigurierbaren Anzahl von Tagen automatisch gelöscht (pro Log individuell einstellbar).
\end{itemize}

\subsection{Laufendes Log}

Das aktuell laufende Log verwendet die Endung \texttt{.log} (z.\,B. \texttt{server.log}).
In der Darstellung \texttt{server.log / ...} bezeichnet der Teil hinter dem Slash die zugehörigen Rollover-Dateien.

Dabei steht \texttt{.i} für eine fortlaufende Nummer und \texttt{yyyy-MM-dd} für das Datum der Logrotation.


\newenvironment{compactalltt}
{\footnotesize\ttfamily
	\setlength{\baselineskip}{0.9\baselineskip}
	\begin{alltt}}
		{\end{alltt}}

\begin{compactalltt}
	log/
	+-- lodoc/
	|   `-- \lodocServerLog
	+-- \mlogisticsLogserverLog
	+-- lomfj/
	|   +-- audit/
	|   |   `-- \lomfjAuditLog
	|   +-- df/
	|   |   `-- \lomfjDfLog
	|   +-- error/
	|   |   `-- \lomfjErrorLog
	|   +-- eventmanager/
	|   |   `-- \lomfjEventmanagerLog
	|   +-- fnas/
	|   |   `-- \lomfjFnasLog
	|   +-- heartbeat/
	|   |   `-- \lomfjHeartbeatLog
	|   +-- incomplete_fnas_pdus/
	|   |   `-- \lomfjIncompleteFnasPdusLog
	|   +-- mfj/
	|   |   `-- \lomfjMfjLog
	|   +-- mfjserver/
	|   |   `-- \lomfjMfjServerLog
	|   +-- pelix/
	|   |   `-- \lomfjPelixLog
	|   `-- uea/
	|       `-- \lomfjUeaLog
	+-- lorouter/
	|   +-- cacheExecutionResult/
	|   |   `-- \lorouterCacheExecutionResultLog
	|   `-- server/
	|       `-- \lorouterServerLog
	+-- mlogistics/
	|   +-- access/
	|   |   `-- \mlogisticsAccessLog
	|   +-- alive/
	|   |   `-- \mlogisticsAliveLog
	|   +-- audit/
	|   |   `-- \mlogisticsAuditLog
	|   +-- commanding/
	|   |   `-- \mlogisticsCommandingLog
	|   +-- commanding_cache/
	|   |   `-- \mlogisticsCommandingCacheLog
	|   +-- commanding_debug/
	|   |   `-- \mlogisticsCommandingDebugLog
	|   +-- cxf/
	|   |   `-- \mlogisticsCxfLog
	|   +-- error/
	|   |   `-- \mlogisticsErrorLog
	|   +-- gdi/
	|   |   `-- \mlogisticsGdiImportLog
	|   +-- locking/
	|   |   `-- \mlogisticsLockingLog
	|   +-- mot/
	|   |   `-- \mlogisticsMotLog
	|   +-- mot_forms/
	|   |   `-- \mlogisticsMotFormsLog
	|   +-- multi_query_servlet/
	|   |   `-- \mlogisticsMultiQueryServletLog
	|   +-- pushd/
	|   |   `-- \mlogisticsPushDLog
	|   +-- query_servlet/
	|   |   `-- \mlogisticsQueryServletLog
	|   +-- query_stream/
	|   |   `-- \mlogisticsQueryStreamLog
	|   +-- scheduler/
	|   |   `-- \mlogisticsSchedulerLog
	|   +-- server/
	|   |   `-- \mlogisticsServerLog
	|   `-- warn/
	|       `-- \mlogisticsWarnLog
	+-- pushd/
	|   +-- error/
	|   |   `-- \pushdErrorLog
	|   `-- server/
	|       `-- \pushdServerLog
	`-- vm-te0001sv0281-1.log
\end{compactalltt}


\subsection{Rollover}

\begin{compactalltt}
	fileserver/log/
	+-- lodoc/
	|   +-- 1970-01-01/
	|   |   `-- \lodocServerLogRollover
	+-- lomfj/
	|   +-- audit/
	|   |   +-- 1970-01-01/
	|   |   |   `-- \lomfjAuditLogRollover
	|   +-- df/
	|   |   `-- 1970-01-01/
	|   |   |   `-- \lomfjDfLogRollover
	|   +-- error/
	|   |   `-- 1970-01-01/
	|   |   |   `-- \lomfjErrorLogRollover
	|   +-- eventmanager/
	|   |   +-- 1970-01-01/
	|   |   |   `-- \lomfjEventmanagerLogRollover
	|   +-- fnas/
	|   |   +-- 1970-01-01/
	|   |   |   `-- \lomfjFnasLogRollover
	|   +-- heartbeat/
	|   |   +-- 1970-01-01/
	|   |   |   `-- \lomfjHeartbeatLogRollover
	|   +-- incomplete_fnas_pdus/
	|   |   +-- 1970-01-01/
	|   |   |   `-- \lomfjIncompleteFnasPdusLogRollover
	|   +-- mfj/
	|   |   +-- 1970-01-01/
	|   |   |   `-- \lomfjMfjLogRollover
	|   +-- mfjserver/
	|   |   +-- 1970-01-01/
	|   |   |   `-- \lomfjMfjServerLogRollover
	|   +-- pelix/
	|   |   +-- 1970-01-01/
	|   |   |   `-- \lomfjPelixLogRollover
	|   `-- uae/
	|       +-- 1970-01-01/
	|           `-- \lomfjUeaLogRollover
	+-- lorouter/
	|   +-- cacheExecutionResult/
	|   |   +-- 1970-01-01/
	|   |   |   `-- \lorouterCacheExecutionResultLogRollover
	|   `-- server/
	|       +-- 1970-01-01/
	|           `-- \lorouterServerLogRollover
	+-- mlogistics/
	|   +-- access/
	|   |   +-- 1970-01-01/
	|   |   |   `-- \mlogisticsAccessLogRollover
	|   +-- alive/
	|   |   +-- 1970-01-01/
	|   |   |   `-- \mlogisticsAliveLogRollover
	|   +-- audit/
	|   |   +-- 1970-01-01/
	|   |   |   `-- \mlogisticsAuditLogRollover
	|   +-- commanding/
	|   |   +-- 1970-01-01/
	|   |   |   `-- \mlogisticsCommandingLogRollover
	|   +-- commanding_cache/
	|   |   +-- 1970-01-01/
	|   |   |   `-- \mlogisticsCommandingCacheLogRollover
	|   +-- commanding_debug/
	|   |   +-- 1970-01-01/
	|   |   |   `-- \mlogisticsCommandingDebugLogRollover
	|   +-- cxf/
	|   |   +-- 1970-01-01/
	|   |   |   `-- \mlogisticsCxfLogRollover
	|   +-- error/
	|   |   +-- 1970-01-01/
	|   |   |   `-- \mlogisticsErrorLogRollover
	|   +-- gdi/
	|   |   +-- 1970-01-01/
	|   |   |   `-- \mlogisticsGdiImportLogRollover
	|   +-- locking/
	|   |   +-- 1970-01-01/
	|   |   |   `-- \mlogisticsLockingLogRollover
	|   +-- mot/
	|   |   +-- 1970-01-01/
	|   |   |   `-- \mlogisticsMotLogRollover
	|   +-- mot_forms/
	|   |   +-- 1970-01-01/
	|   |   |   `-- \mlogisticsMotFormsLogRollover
	|   +-- multi_query_servlet/
	|   |   +-- 1970-01-01/
	|   |   |   `-- \mlogisticsMultiQueryServletLogRollover
	|   +-- pushd/
	|   |   +-- 1970-01-01/
	|   |   |   `-- \mlogisticsPushDLogRollover
	|   +-- query_servlet/
	|   |   +-- 1970-01-01/
	|   |   |   `-- \mlogisticsQueryServletLogRollover
	|   +-- query_stream/
	|   |   +-- 1970-01-01/
	|   |   |   `-- \mlogisticsQueryStreamLogRollover
	|   +-- scheduler/
	|   |   +-- 1970-01-01/
	|   |   |   `-- \mlogisticsSchedulerLogRollover
	|   +-- server/
	|   |   +-- 1970-01-01/
	|   |   |   `-- \mlogisticsServerLogRollover
	|   `-- warn/
	|       +-- 1970-01-01/
	|           `-- \mlogisticsWarnLogRollover
	+-- pushd/
	|   +-- error/
	|   |   +-- 1970-01-01/
	|   |   |   `-- \pushdErrorLogRollover
	|   `-- server/
	|       +-- 1970-01-01/
	|           `-- \pushdServerLogRollover
\end{compactalltt}
